\documentclass[10pt,a4paper]{amsart}
\usepackage[margin=19mm]{geometry}
\usepackage{amsmath,amssymb,mathtools}
\usepackage[scr=rsfs,cal=euler]{mathalfa}
\usepackage{xfrac}
\usepackage[unicode,hidelinks,bookmarks=false]{hyperref}
\newcommand{\ie}{i.e.\ }
\newcommand{\cf}{cf.\ }
\newcommand{\resp}{respectively\ }
\newcommand{\Subsection}{\subsection}
\providecommand{\nobreakdash}{\nobreak\hbox{-}}
\setlength{\emergencystretch}{2em}
\multlinegap0pt
\usepackage{mathtools}
\usepackage{amssymb}
\usepackage{xcolor}
\newtheorem{lemma}{Lemma}
\title{Erd\H{o}s's diameter conjecture for separated distances fails in high dimensions}
\author{Boon Suan Ho}
\date{Source: arXiv:2604.15305v1 (CC BY 4.0)}
\begin{document}
\maketitle

\noindent\emph{Source and licence.} Erd\H{o}s's diameter conjecture for separated distances fails in high dimensions. Version arXiv:2604.15305v1 (CC BY 4.0), distributed by arXiv under the Creative Commons Attribution 4.0 International licence, http://creativecommons.org/licenses/by/4.0/, as recorded in the arXiv licence metadata for this exact version and shown on the abstract page https://arxiv.org/abs/2604.15305v1. Full licence notice: \texttt{licenses/arxiv-2604.15305-CC-BY-4.0.txt}.\par\smallskip

\noindent\textit{Editorial scope.} Counted unit: the article's lemma lem:monotone-concave (source0.tex lines 263--265). The article's complete original proof of it is delivered with it (source0.tex lines 267--310, one \texttt{proof} environment of 44 lines). No mathematics is written, repaired or re-derived: the statement and the proof are the article's own lines, in the article's own notation, and no retained line is substituted or re-flowed.  No further source-local context is needed: every cross-reference of every retained line resolves inside the two delivered spans.  Nothing was omitted from any retained span, so the package carries no omission record.  No retained line contains a citation, so the wrapper carries no bibliography and no bibliography entry.  No retained line contains a figure environment, an image-inclusion command or drawing code, so no image is delivered and the package declares no asset dependency.  Editorial formatting only, in addition: an AMS \texttt{amsart} preamble that restates the article's own declarations and macros used by the retained lines, namely the environments \texttt{lemma} on separate counters, together with the standard packages those lines need.  The retained environments print in this document as Lemma 1 in order of appearance, whereas the article prints the same items with the numbering of its own full document.  The complete native source of the article as distributed in the arXiv e-print for this exact version is bundled unchanged as \texttt{sources/198598/source0.tex}.
\begin{lemma}\label{lem:monotone-concave}
The function $G$ is strictly increasing and strictly concave on $(0,\pi)$.
\end{lemma}

\begin{proof}
Set $a\coloneqq\pi/4$, so that
\[
H(\theta)=a\theta-\varepsilon(1-\cos\theta).
\]
We have
\[
H'(\theta)=a-\varepsilon\sin\theta\ge a-\varepsilon=\frac{\pi}{4}-\frac18>0,
\]
so $H$, and hence $G$, is strictly increasing on $(0,\pi)$.

For concavity we use
\begin{equation}\label{eq:G-second}
G''(\theta)=\frac{2H(\theta)H''(\theta)-H'(\theta)^2}{4H(\theta)^{3/2}}.
\end{equation}
Since $H''(\theta)=-\varepsilon\cos\theta$, a short calculation gives
\begin{equation}\label{eq:numerator-factorized}
2H(\theta)H''(\theta)-H'(\theta)^2
=-a^2+2a\varepsilon\bigl(\sin\theta-\theta\cos\theta\bigr)
-\varepsilon^2(1-\cos\theta)^2.
\end{equation}
Now set
\[
\phi(\theta)\coloneqq\sin\theta-\theta\cos\theta.
\]
Then
\[
\phi'(\theta)=\theta\sin\theta>0
\qquad (0<\theta<\pi),
\]
so $\phi$ is strictly increasing on $(0,\pi)$ and therefore
\[
0<\phi(\theta)<\phi(\pi)=\pi.
\]
Because $\varepsilon=1/8=a/(2\pi)$, equation
\eqref{eq:numerator-factorized} implies
\[
2H(\theta)H''(\theta)-H'(\theta)^2
\le -a^2+2a\varepsilon\pi-\varepsilon^2(1-\cos\theta)^2
=-\varepsilon^2(1-\cos\theta)^2<0
\]
for every $0<\theta<\pi$. By \eqref{eq:G-second}, this proves that $G$ is
strictly concave.
\end{proof}

\end{document}
